
\section{Fundamentos da cinemática}
\label{sec:fundamentos}

A cinemática de um manipulador robótico permite estabelecer a relação entre as variáveis articulares e a posição e orientação do \textit{end effector}. No caso do myCobot~280, constituído por seis articulações rotativas, esta análise divide-se em dois problemas fundamentais: a cinemática direta e a cinemática inversa.

A cinemática direta permite determinar a posição e orientação de um manipulador a partir dos ângulos das suas articulações. Por outro lado, a cinemática inversa permite determinar os ângulos articulares necessários para alcançar uma posição pretendida, considerando determinadas restrições de orientação.

